The advent of large language models has substantially lowered the barrier to KG
construction at scale, and the popularity of this approach has led to a proliferation of tools and frameworks for LLM-augmented KG construction.
\citet{pan2024llmkg} survey this landscape, classifying LLM-augmented construction methods
by the role of the LLM --- entity discovery, relation extraction, end-to-end generation, or
knowledge distillation. Among open-source frameworks, LlamaIndex, LangChain, and Neo4j are the dominant providers.
LlamaIndex~\citep{liu2022llamaindex} offers three extraction strategies for property graph
construction~\citep{llamaindex2024propertygraph}: a schema-free extractor in which the LLM
assigns entity and relation types freely (\texttt{SimpleLLMPathExtractor}); a strict schema
extractor that enforces a user-defined list of allowed types
(\texttt{SchemaLLMPathExtractor}); and a guided extractor that accepts an initial type list
as non-binding hints while permitting the LLM to introduce novel types
(\texttt{DynamicLLMPathExtractor}).
LangChain's \texttt{LLMGraphTransformer}~\citep{chase2022langchain,
bratanic2024llmgraphtransformer} follows a similar pattern.
In both modes users may supply \texttt{allowed\_nodes} and \texttt{allowed\_relationships}
as schema constraints, with a strict flag that discards extractions violating the schema.
The Neo4j LLM Knowledge Graph Builder~\citep{neo4j2024llmgraphbuilder,
neo4jdocs2024kgbuilder} offers three schema modes: a user-supplied schema of node types,
relationship types, and valid source--relation--target triples; an auto-extracted schema
inferred from the input text in a single upfront LLM call and held fixed thereafter; and a
free mode with no schema.

These tools have seen wide adoption but share a fundamental design choice: the ontology
--- node types, relation types, and constraints --- is specified arbitrarily and remains fixed for
the duration of construction or even sometimes is built ad hoc by the LLM without an explicit schema.

Several academic approaches try to bridge triple extraction quality and task specification.
\citet{mo2025kggen} propose KGGen, which extracts KGs from plain text using language models
with improved entity and relation normalisation: an iterative clustering step collapses the
near-duplicate relation types that open extraction produces, yielding denser and more
connected graphs.
\citet{mynarz2023testdriven} propose a test-driven approach in which competency questions,
formalised as SPARQL queries, drive extraction --- providing a
form of task specification without a quantitative feedback loop.

\citet{li2025llmontologyreview} provide the most comprehensive survey to date of LLMs
across the ontology engineering lifecycle, reviewing 30 publications and 41 experiments
spanning requirements specification, conceptualisation, encoding, ontology matching, and
maintenance.
The distribution of effort is stark: ontology implementation accounts for 63.4\% of studied
experiments, while ontology revision and maintenance has been
addressed by a single study (2.4\%).
The survey names the consequences directly: ``support across the ontology lifecycle is
uneven, with maintenance particularly underexplored,'' and finds that ``evaluation practices
are fragmented, as existing quantitative metrics fail to fully capture
performance''~\citep{li2025llmontologyreview}; across all reviewed studies, precision,
recall, and F1 are applied against ontology benchmarks such as the Ontology Alignment
Evaluation Initiative (OAEI), not against a downstream application.

The survey closes by identifying ontology revision and maintenance as an open direction: 
``Future research directions may include developing hybrid learning
pipelines that allow LLMs to continuously integrate domain updates \ldots\ thereby supporting
ongoing validation and refinement based on domain-specific
feedback''~\citep{li2025llmontologyreview}.

The evaluation of knowledge graphs has historically been conducted through intrinsic,
structural quality metrics. \citet{paulheim2017kgquality} survey methods for KG completion
and error detection, evaluating them against structural criteria such as precision and recall
on missing type and relation assertions. More recently,
\citet{paulheim2024structuralquality} draw an explicit distinction between this intrinsic
mode of evaluation and task-based evaluation --- assessing a KG by its utility on
a concrete downstream application. This distinction is central to the motivation of the
present work: a graph may score well on structural criteria while failing to support a
specific predictive task.

Several frameworks try to standardise task-based KG evaluation. KGBench~\cite{bloem2021kgbench}
provides a suite of relational and multimodal machine learning tasks benchmarked on
curated KG datasets.
KGrEaT~\cite{heist2023kgreat} enables systematic comparison of KG variants across
classification, regression, and recommendation tasks; the framework explicitly supports
comparison of ``different versions of a KG during its life cycle,'' precisely the kind of
comparison our feedback loop performs automatically. LLM-KG-Bench~\cite{aksw2024llmkgbench}
focuses specifically on the quality of LLM-generated graphs. Together, these frameworks
confirm that task utility cannot be inferred from structural properties alone --- yet none
of them feeds performance signals back into the construction process itself.

The closest prior work to our approach concerns task-aware graph construction.
\citet{tsang2025autographr1} introduce AutoGraph-R1, the first framework to directly
optimise KG construction for task performance using reinforcement learning: an LLM
constructor is trained as a policy whose reward derives from the graph's functional utility
in a downstream RAG retrieval pipeline.
AutoGraph-R1 shares with our approach the core principle of aligning construction with a
downstream objective; the critical differences are that its reward derives from retrieval
performance in a RAG pipeline rather than a node-level classifier, the RL policy optimises
extraction density and triple relevance within an open-schema setting rather than selecting
among ontology configurations, and there is no mechanism for explicit schema evolution.

\citet{cucumides2025augraph} propose constructing graphs from tabular and relational data
to maximise GNN performance on a target task, tailoring graph topology through iterative
process where each attribute's predictive
relevance is assessed. However, their setting involves structured tabular input rather than unstructured
text, and the construction is not driven by an evolving ontology.

The financial domain chosen for this study provides a natural testbed for task-conditioned KG construction.
Financial news significantly influences market dynamics~\citep{dong2024fnspid}, and a long
line of work has explored how to represent its informational content as structured graphs
for downstream prediction.
\citet{ding2015stockpred} demonstrate that event-driven representations of financial
news --- structured as subject--predicate--object triples capturing who did what to whom
--- yield strong signals for short- and medium-term stock price movement when combined with deep neural
networks.
Subsequent work confirms and extends this finding: \citet{ding2016kgeventstock} show that
enriching event representations with knowledge graph embeddings further improves prediction
accuracy, and attention-based models over news sequences~\citep{hu2018chaoticwhispers} and
financial social media text~\citep{xu2018stocktweets} demonstrate the breadth of the
predictive signal across sources and markets.

\citet{arun2025finreflectkg} introduce FinReflectKG, a pipeline for constructing financial
KGs from SEC 10-K filings using a pre-configured schema of entity and relationship types
specific to financial regulatory documents.
An LLM-based reflection component iterates to correct extraction errors within this fixed
schema until no issues are detected or a step budget is exhausted.

Taken together, the literature reveals a consistent gap.
Evaluation frameworks confirm that KG utility is task-dependent and cannot be inferred
from structural properties alone~\citep{heist2023kgreat, bloem2021kgbench}.
Task-aware construction methods demonstrate that graph structure can be aligned with
downstream objectives~\citep{cucumides2025augraph, tsang2025autographr1}, and
many researchers assume that event-driven representations of financial news carry
strong predictive signals for market classification~\citep{ding2015stockpred, ding2016kgeventstock,
hu2018chaoticwhispers, xu2018stocktweets}.
Yet no existing work closes the loop: temporal KG methods evolve facts within a fixed
schema~\citep{trivedi2017knowevolve}; task-aware methods adapt topology or extraction reward
without revising the ontology~\citep{cucumides2025augraph, tsang2025autographr1};
financial KG systems such as FinReflectKG~\citep{arun2025finreflectkg} iterate within a
pre-configured schema; and ontology engineering surveys identify schema maintenance as the
most underexplored stage of the lifecycle~\citep{li2025llmontologyreview}.
No existing method treats the KG ontology --- the choice of node types, relation types,
and constraints --- as a hyperparameter to be optimised through a closed feedback loop
between a downstream classifier and the construction process.
The present paper addresses this gap directly.